\chapter{Vision 2030: The Phase Transition Realised}
\label{ch:vision}

\epigraph{Prediction is very difficult, especially about the future.}{attributed
to Niels Bohr}

This chapter discharges contribution~C6: a 2030 outlook. It is the most
speculative chapter, and it is marked accordingly. Forecasts from analysts and
market researchers are reproduced as \forecast; the thesis's own extrapolations
are derived from the model of \cref{ch:framework,ch:math} and are likewise bounded.
\Cref{sec:vis-forecasts} reviews external forecasts; \cref{sec:vis-expansion}
applies the phase-transition model; \cref{sec:vis-roadmap} gives a staged roadmap;
and \cref{sec:vis-governance} treats governance and risk.

\section{External Forecasts}
\label{sec:vis-forecasts}

The analyst consensus for the late 2020s is one of rapid agentic adoption paired
with high failure rates. Gartner projects that by 2028 one-third of enterprise
software applications will include agentic AI, up from less than one per cent in
2024, and that at least fifteen per cent of day-to-day work decisions will be made
autonomously \citep{gartner-agentic-2025}~[\forecast]; that forty per cent of
enterprise applications will feature task-specific agents by 2026
\citep{gartner-taskagents-2025}~[\forecast]; and, in the same breath, that over
forty per cent of agentic-AI projects will be cancelled by the end of 2027, citing
``inadequate risk controls'' among the causes
\citep{gartner-agentic-2025}~[\forecast]. Most consequentially for this thesis,
Gartner forecasts that ``guardian agent'' technologies will capture ten to fifteen
per cent of the agentic-AI market by 2030 \citep{gartner-guardian-2025}~[\forecast].
Market sizings vary widely by definition: one estimate puts the AI-agents market at
USD~7.84\,billion in 2025 rising to USD~52.62\,billion by 2030 (CAGR ${\approx}46\%$)
\citep{marketsandmarkets-aiagents}~[\forecast], another the enterprise segment at
USD~2.58\,billion rising to USD~24.5\,billion \citep{grandview-agentic}~[\forecast].
The spread (a factor of two at the 2030 horizon) is itself the finding: these are
indicative ranges, not measurements.

Two of these forecasts bear directly on the argument. The ``inadequate risk
controls'' cancellation driver is precisely the deficiency P2 of
\cref{sec:problem}; and the ``guardian agent'' category is, in this thesis's
vocabulary, the market's name for a law-state runtime. The external forecast that a
distinct stratum of oversight agents will command a tenth of the market by 2030 is
independent, if speculative, support for \cref{hyp:substrate}'s claim that
conformance is a stratum in its own right.

\section{The Expansion Argument}
\label{sec:vis-expansion}

What does the phase-transition model predict, taken seriously and stated honestly?
The Decoupling Theorem (\cref{thm:decoupling}) and its corollary
(\cref{cor:superlinear}) establish that, once domains, capabilities, and agents are
each integrated once against the fused substrate, the reachable surface scales as
the product $\Theta(n^3)$ against linear integration cost $\Theta(n)$. The Landau
analysis (\cref{prop:landau-jump}) establishes that, with a network-effect term,
the adoption order parameter jumps discontinuously as standardisation crosses a
threshold. Composing the two: \emph{if} adoption crosses its threshold in the late
2020s---of which the institutional consolidation under the Linux Foundation and the
Agentic AI Foundation (\cref{sec:vis-governance}) is a leading indicator---then the
addressable surface of machine-mediated language expands not incrementally but by
orders of magnitude, the digital counterpart of the ${\sim}1{,}600$--$1{,}700\times$
expansion that names this thesis.

This is the steam claim, and three honesty constraints bound it. First, it is a
\forecast: the model explains a mechanism, it does not measure $\Phi$. Second, by
\cref{prin:plateau}, the present period of intense standardisation effort with
modest visible capability is consistent with an \emph{approaching} transition and
with one that never arrives; the same plateau precedes both. Third, the
order-of-magnitude figure is borrowed from physics as an emblem of \emph{character}
(discontinuous, large), not as a measured digital quantity. With those constraints,
the prediction is: the convergence of LSP, MCP, and A2A under a conformance stratum
is of the order type of a first-order phase transition, and the late 2020s are
when the threshold is most plausibly crossed.

\section{A Roadmap to 2030}
\label{sec:vis-roadmap}

\Cref{tab:roadmap} stages the transition with bounded entry criteria, so that the
roadmap is falsifiable rather than aspirational: each stage names what would count
as having entered it.

\begin{table}[t]
\centering
\caption[A staged roadmap to 2030]{A staged roadmap. Each stage's status is the
bounded status of its entry criterion as of submission; later stages are
\openstat\ by construction.}
\label{tab:roadmap}
\small
\begin{tabularx}{\textwidth}{@{}l Y l@{}}
\toprule
\textbf{Stage} & \textbf{Entry criterion} & \textbf{Status} \\
\midrule
S1 (now) & Generalised LSP servers exist for non-code languages; agent traces emitted as OCEL & \candidate \\
S2 & Read-only conformance exposed as MCP tools/resources; receipts cross tool boundaries & \candidate \\
S3 & Conformance authorities discoverable over A2A; verdicts portable across agents & \openstat \\
S4 & Cross-domain law axes standardised under a neutral foundation & \openstat \\
S5 (2030) & Conformance-for-all-language is a default stratum of the substrate & \openstat \\
\bottomrule
\end{tabularx}
\end{table}

Stage~S1 is already underway: language servers for natural language and data
formats exist \citep{ltex,yamlls}, and the artifact of \cref{ch:artifact} emits its
own agent conduct as OCEL. Stage~S2 is foreshadowed by the workflow-engine and
capability-binding preprints of \cref{sec:recent} \citep{parmar26,zhou26gov}.
Stage~S3 awaits A2A-advertised conformance capabilities, of which the
identity-and-delegation work across MCP and A2A is a precursor
\citep{prakash26aip}. Stages~S4 and~S5 are \openstat: they require cross-domain law
standardisation that does not yet exist, and the thesis does not claim it will.

\section{Governance, Standardisation, and Risk}
\label{sec:vis-governance}

The latent heat of standardisation (\cref{def:latent}) is being paid through
institutions. A2A entered the Linux Foundation in June 2025
\citep{a2a-lf-donation,lf-a2a-press}; MCP entered the Agentic AI Foundation under
the same umbrella in December 2025 \citep{mcp-aaif-2025,lf-aaif-press}. Neutral
governance is the social mechanism by which a protocol's adoption threshold is
crossed, and its arrival for two of the three strata within six months is the
clearest leading indicator that the transition's precondition is being met.

The risks are equally concrete, and conformance speaks to several. The cancellation
forecast attributes failure partly to weak risk controls
\citep{gartner-agentic-2025}; the security literature documents prompt injection
and tool poisoning against MCP \citep{gaire25,huang26threat,maloyan26}; the
reliability literature finds that accuracy masks unreliability \citep{rabanser26};
and the oversight literature shows that human oversight has finite capacity, so
that adding escalation can \emph{reduce} safety past a point \citep{turan26}. A
law-state runtime is a partial answer to the first three---a standing, three-valued,
receipt-bound risk control---and the ``guardian agent'' market
\citep{gartner-guardian-2025} is where it would live. But it is no answer to the
fourth: an oversight stratum that itself demands human attention inherits the
capacity limit, and ``proof-of-guardrail'' can be trusted only as far as its own
analysis allows \citep{jin26guardrail}. The honest 2030 vision is therefore
double-edged: the same fusion that expands reach also concentrates a new dependency
on the conformance stratum, whose own trustworthiness, scalability, and capture
resistance are \openstat. Whether the guardian guards itself is the question the
next chapter's threats-to-validity must hold open.
